\documentclass[a4paper,12pt]{article}

\usepackage{cmap}
\usepackage{mathtext}
\usepackage[T2A]{fontenc}
\usepackage[utf8]{inputenc}
\usepackage[english,russian]{babel}
\usepackage{indentfirst}

\usepackage{extsizes}
\usepackage{geometry}
\geometry{top=25mm}
\geometry{bottom=35mm}
\geometry{left=35mm}
\geometry{right=20mm}

\usepackage{setspace}
\usepackage{lastpage}
\usepackage{soul}
\usepackage{multicol}
\usepackage{lipsum}

\usepackage{hyperref}
\usepackage[usenames,dvipsnames,svgnames,table,rgb]{xcolor}
\hypersetup{
  unicode=true,
  pdftitle={Заголовок},
  pdfauthor={Автор},
  pdfsubject={Тема},
  pdfcreator={Создатель},
  pdfproducer={Производитель},
  pdfkeywords={keyword1} {key2} {key3},
  colorlinks=true,
  linkcolor=black,
  citecolor=green,
  filecolor=magenta,
  urlcolor=blue
}

\usepackage{amsfonts,amssymb,amsthm,mathtools}
\usepackage{amsmath}
\usepackage{icomma}

\usepackage{euscript}
\usepackage{mathrsfs}

\newcommand*{\hm}[1]{#1\nobreak\discretionary{}
{\hbox{$\mathsurround=0pt #1$}}{}}

\usepackage{graphicx}
\graphicspath{{imgs/}}
\setlength\fboxsep{3pt}
\setlength\fboxrule{1pt}
\usepackage{wrapfig}

\author{ФИО}
\title{\LaTeX. Вставка и настройка рисунков}
\date{\today}

\begin{document}

\maketitle

\section{Форматирование рисунков}

\noindent\includegraphics{picture}

\bigskip

\noindent Увеличение в 3 раза:\\
\includegraphics[scale=3]{picture}

\bigskip

\noindent Уменьшение в 2 раза:\\
\includegraphics[scale=0.5]{picture}

\bigskip

\noindent Ширина 2 см, высота 5 см (с сохранением пропорций):\\
\includegraphics[width=2cm,height=5cm,keepaspectratio]{picture}

\bigskip

\noindent Рисунок в рамке (отступ 10 пт, толщина 2 пт):\\
{\setlength\fboxsep{10pt}\setlength\fboxrule{2pt}\fbox{\includegraphics[scale=0.5]{picture}}}

\section{Нумерация рисунков}

См. рисунок~\ref{fig:numbered} на странице~\pageref{fig:numbered}.

\begin{figure}[htbp]
  \centering
  \includegraphics[scale=0.5]{picture}
  \caption{Пример нумерованного рисунка}
  \label{fig:numbered}
\end{figure}

\section{Обтекаемый рисунок}

\lipsum[1-2]

\begin{wrapfigure}{r}{0.35\textwidth}
  \centering
  \caption{Обтекаемый рисунок}
  \includegraphics[width=0.3\textwidth]{picture}
\end{wrapfigure}

\lipsum[3-4]

\section{Два рисунка под одним названием}

\begin{figure}[htbp]
  \centering
  \begin{minipage}{0.45\textwidth}
  \center{\includegraphics[width=\textwidth]{picture} \\ а)}
\end{minipage}\hfill
\begin{minipage}{0.45\textwidth}
\center{ \includegraphics[width=\textwidth]{picture} \\ б) }
\end{minipage}
\caption{Два рисунка под одним названием}
\label{fig:two}
\end{figure}

\listoffigures

\end{document}